% ============================================================
\chapter{Mesures}
\label{ch:mesures}
% ============================================================

Maintenant que nous disposons des $\sigma$-alg\`ebres --- les collections
d'ensembles auxquels nous souhaitons attribuer une taille --- il est temps
de d\'efinir cette \og taille\fg{} elle-m\^eme. C'est l'objet de la notion de
\emph{mesure}, qui g\'en\'eralise simultan\'ement la longueur, l'aire, le
volume et la probabilit\'e. Cette unification, due \`a Lebesgue et
perfectionn\'ee par Carath\'eodory, est l'une des grandes r\'eussites des
math\'ematiques modernes.

% ------------------------------------------------------------
\section{D\'efinition et premi\`eres propri\'et\'es}
% ------------------------------------------------------------

\begin{definition}[Mesure]
Soit $(X, \mathcal{A})$ un espace mesurable. Une \emph{mesure} sur
$(X, \mathcal{A})$ est une application $\mu : \mathcal{A} \to [0, +\infty]$
v\'erifiant~:
\begin{enumerate}
    \item[(i)] $\mu(\varnothing) = 0$\,;
    \item[(ii)] ($\sigma$-additivit\'e) pour toute suite $(A_n)_{n \in \N}$
    d'\'el\'ements de $\mathcal{A}$ \emph{deux \`a deux disjoints}~:
    \[
    \mu\!\left(\bigsqcup_{n=0}^{\infty} A_n\right) = \sum_{n=0}^{\infty} \mu(A_n).
    \]
\end{enumerate}
Le triplet $(X, \mathcal{A}, \mu)$ est appel\'e \emph{espace mesur\'e}.
\end{definition}

\begin{definition}[Types de mesures]
Soit $(X, \mathcal{A}, \mu)$ un espace mesur\'e.
\begin{itemize}
    \item $\mu$ est \emph{finie} si $\mu(X) < \infty$.
    \item $\mu$ est une \emph{mesure de probabilit\'e} si $\mu(X) = 1$.
    \item $\mu$ est \emph{$\sigma$-finie} s'il existe une suite $(X_n)_{n \in \N}
    \subset \mathcal{A}$ avec $X = \bigcup_n X_n$ et $\mu(X_n) < \infty$ pour
    tout $n$.
    \item $\mu$ est \emph{compl\`ete} si pour tout $N \in \mathcal{A}$ avec
    $\mu(N) = 0$ et tout $A \subset N$, on a $A \in \mathcal{A}$.
\end{itemize}
\end{definition}

\begin{proposition}[Monotonie et sous-additivit\'e]
\label{prop:monotonie}
Soit $(X, \mathcal{A}, \mu)$ un espace mesur\'e.
\begin{enumerate}
    \item \textbf{Monotonie}~: si $A \subset B$ avec $A, B \in \mathcal{A}$, alors
    $\mu(A) \leq \mu(B)$.
    \item \textbf{Diff\'erence}~: si $A \subset B$ et $\mu(A) < \infty$, alors
    $\mu(B \setminus A) = \mu(B) - \mu(A)$.
    \item \textbf{Sous-additivit\'e d\'enombrable}~: pour toute suite
    $(A_n) \subset \mathcal{A}$,
    $\mu\!\left(\bigcup_n A_n\right) \leq \sum_n \mu(A_n)$.
\end{enumerate}
\end{proposition}

\begin{proof}
(1) On \'ecrit $B = A \sqcup (B \setminus A)$, donc
$\mu(B) = \mu(A) + \mu(B \setminus A) \geq \mu(A)$.

(2) D\'ecoule directement de la d\'ecomposition $B = A \sqcup (B \setminus A)$
et de la finitude de $\mu(A)$.

(3) Posons $B_0 = A_0$ et $B_n = A_n \setminus \bigcup_{k=0}^{n-1} A_k$ pour
$n \geq 1$. Alors $(B_n)$ est une suite disjointe avec $\bigcup_n B_n =
\bigcup_n A_n$ et $B_n \subset A_n$. Donc
\[
\mu\!\left(\bigcup_n A_n\right) = \mu\!\left(\bigsqcup_n B_n\right)
= \sum_n \mu(B_n) \leq \sum_n \mu(A_n). \qedhere
\]
\end{proof}

\section{Continuit\'e de la mesure}

\begin{theoreme}[Continuit\'e monotone]
\label{thm:continuite_monotone}
Soit $(X, \mathcal{A}, \mu)$ un espace mesur\'e.
\begin{enumerate}
    \item \textbf{Continuit\'e croissante}~: si $A_n \uparrow A$ (suite croissante),
    alors $\mu(A_n) \uparrow \mu(A)$.
    \item \textbf{Continuit\'e d\'ecroissante}~: si $A_n \downarrow A$ (suite
    d\'ecroissante) et $\mu(A_1) < \infty$, alors $\mu(A_n) \downarrow \mu(A)$.
\end{enumerate}
\end{theoreme}

\begin{proof}
(1) Posons $B_0 = A_0$ et $B_n = A_n \setminus A_{n-1}$ pour $n \geq 1$. Alors
$(B_n)$ est disjointe et $\bigsqcup_{k=0}^n B_k = A_n$, $\bigsqcup_{k=0}^\infty
B_k = A$. Par $\sigma$-additivit\'e~:
\[
\mu(A) = \sum_{k=0}^{\infty} \mu(B_k) = \lim_{n \to \infty} \sum_{k=0}^n \mu(B_k)
= \lim_{n \to \infty} \mu(A_n).
\]

(2) Posons $C_n = A_1 \setminus A_n$. Alors $C_n \uparrow A_1 \setminus A$.
Par (1), $\mu(C_n) \uparrow \mu(A_1 \setminus A)$. Comme $\mu(A_1) < \infty$~:
$\mu(C_n) = \mu(A_1) - \mu(A_n)$ et $\mu(A_1 \setminus A) = \mu(A_1) - \mu(A)$,
d'o\`u $\mu(A_n) \downarrow \mu(A)$.
\end{proof}

\begin{attention}[title=\textbf{Condition de finitude}]
L'hypoth\`ese $\mu(A_1) < \infty$ dans (2) est essentielle. Contre-exemple~:
$A_n = [n, +\infty)$ avec la mesure de Lebesgue. On a $A_n \downarrow \varnothing$
mais $\mu(A_n) = +\infty$ pour tout $n$.
\end{attention}

\begin{lemme}[Lemme de Borel--Cantelli]
\label{lem:borel_cantelli}
Soit $(X, \mathcal{A}, \mu)$ un espace mesur\'e et $(A_n) \subset \mathcal{A}$. Si
$\sum_{n=0}^{\infty} \mu(A_n) < \infty$, alors
\[
\mu\!\left(\limsup_{n \to \infty} A_n\right) = 0,
\]
o\`u $\limsup_n A_n = \bigcap_{n=0}^{\infty} \bigcup_{k=n}^{\infty} A_k$ est
l'ensemble des points appartenant \`a une infinit\'e de $A_n$.
\end{lemme}

\begin{proof}
Posons $B_n = \bigcup_{k=n}^{\infty} A_k$. Alors $B_n \downarrow \limsup_n A_n$ et
\[
\mu(B_n) \leq \sum_{k=n}^{\infty} \mu(A_k) \to 0 \quad (n \to \infty)
\]
car le reste d'une s\'erie convergente tend vers z\'ero. Comme $\mu(B_0) \leq
\sum_k \mu(A_k) < \infty$, la continuit\'e d\'ecroissante donne
$\mu(\limsup_n A_n) = \lim_n \mu(B_n) = 0$.
\end{proof}

\section{Exemples fondamentaux}

\begin{exemple}[Mesure de comptage]
Soit $X$ un ensemble quelconque et $\mathcal{A} = \mathcal{P}(X)$. La \emph{mesure
de comptage} est d\'efinie par~:
\[
\mu(A) = \begin{cases} \abs{A} & \text{si } A \text{ est fini,} \\
+\infty & \text{sinon.} \end{cases}
\]
C'est une mesure sur $(X, \mathcal{P}(X))$, $\sigma$-finie si et seulement si $X$
est d\'enombrable.
\end{exemple}

\begin{exemple}[Masse de Dirac]
Soit $x_0 \in X$. La \emph{masse de Dirac} en $x_0$ est la mesure de probabilit\'e~:
\[
\delta_{x_0}(A) = \begin{cases} 1 & \text{si } x_0 \in A, \\ 0 & \text{si }
x_0 \notin A. \end{cases}
\]
\end{exemple}

\begin{exemple}[Combinaisons lin\'eaires]
Si $\mu_1, \mu_2$ sont des mesures sur $(X, \mathcal{A})$ et $\alpha, \beta \geq 0$,
alors $\alpha \mu_1 + \beta \mu_2$ est une mesure. Plus g\'en\'eralement, si
$(x_n)_{n \in \N} \subset X$ et $(a_n)_{n \in \N} \subset [0, +\infty)$, alors
$\mu = \sum_n a_n \delta_{x_n}$ est une mesure.
\end{exemple}

\section{Pr\'e-mesures et extension}

\begin{definition}[Pr\'e-mesure]
Soit $\mathcal{A}_0$ une alg\`ebre sur $X$. Une application
$\mu_0 : \mathcal{A}_0 \to [0, +\infty]$ est une \emph{pr\'e-mesure} si~:
\begin{enumerate}
    \item[(i)] $\mu_0(\varnothing) = 0$\,;
    \item[(ii)] pour toute suite $(A_n) \subset \mathcal{A}_0$ disjointe telle que
    $\bigsqcup_n A_n \in \mathcal{A}_0$, on a
    $\mu_0\!\left(\bigsqcup_n A_n\right) = \sum_n \mu_0(A_n)$.
\end{enumerate}
\end{definition}

\begin{remarque}
La diff\'erence avec une mesure est que $\mathcal{A}_0$ est seulement une alg\`ebre,
et la $\sigma$-additivit\'e n'est exig\'ee que lorsque la r\'eunion tombe dans
$\mathcal{A}_0$.
\end{remarque}

\begin{proposition}[Crit\`ere de $\sigma$-additivit\'e]
\label{prop:critere_sigma_add}
Soit $\mu_0$ une application finiment additive sur une alg\`ebre $\mathcal{A}_0$
avec $\mu_0(\varnothing) = 0$. Les conditions suivantes sont \'equivalentes~:
\begin{enumerate}
    \item $\mu_0$ est une pr\'e-mesure ($\sigma$-additive).
    \item $\mu_0$ est continue par en bas~: si $(A_n) \subset \mathcal{A}_0$,
    $A_n \uparrow A \in \mathcal{A}_0$, alors $\mu_0(A_n) \uparrow \mu_0(A)$.
    \item $\mu_0$ est continue en $\varnothing$~: si $(A_n) \subset \mathcal{A}_0$,
    $A_n \downarrow \varnothing$, alors $\mu_0(A_n) \downarrow 0$.
\end{enumerate}
\end{proposition}

\begin{proof}
$(1) \Rightarrow (2)$~: m\^eme preuve que le th\'eor\`eme~\ref{thm:continuite_monotone}(1).

$(2) \Rightarrow (3)$~: si $A_n \downarrow \varnothing$ et $\mu_0(A_1) < \infty$,
on applique la m\^eme technique que~\ref{thm:continuite_monotone}(2). Si
$\mu_0(A_1) = +\infty$, le r\'esultat est trivial.

$(3) \Rightarrow (1)$~: soit $(A_n)$ disjointe avec $A = \bigsqcup_n A_n \in
\mathcal{A}_0$. Posons $B_N = A \setminus \bigsqcup_{n=0}^{N} A_n$. Par additivit\'e
finie, $\mu_0(A) = \sum_{n=0}^{N} \mu_0(A_n) + \mu_0(B_N)$. Comme $B_N \downarrow
\varnothing$, on a $\mu_0(B_N) \to 0$, d'o\`u $\mu_0(A) = \sum_n \mu_0(A_n)$.
\end{proof}

\section{Th\'eor\`eme d'extension}

\begin{theoreme}[Extension de Carath\'eodory --- \'enonc\'e pr\'eliminaire]
\label{thm:extension_preliminaire}
Soit $\mu_0$ une pr\'e-mesure sur une alg\`ebre $\mathcal{A}_0$. Alors $\mu_0$
s'\'etend en une mesure $\mu$ sur $\sigma(\mathcal{A}_0)$. Si $\mu_0$ est
$\sigma$-finie, cette extension est unique.
\end{theoreme}

La preuve compl\`ete sera donn\'ee au chapitre~\ref{ch:caratheodory}, mais nous
esquissons ici la construction.

\begin{formules}[title=\textbf{Construction de la mesure ext\'erieure}]
On d\'efinit la \emph{mesure ext\'erieure} associ\'ee \`a $\mu_0$ par~:
\[
\mu^*(A) = \inf\left\{\sum_{n=0}^{\infty} \mu_0(A_n) :
(A_n) \subset \mathcal{A}_0,\, A \subset \bigcup_{n=0}^{\infty} A_n\right\}
\]
pour tout $A \subset X$. Puis on restreint $\mu^*$ aux ensembles
\emph{Carath\'eodory-mesurables}.
\end{formules}

\section{Compl\'etion d'une mesure}

\begin{definition}[Compl\'etion]
Soit $(X, \mathcal{A}, \mu)$ un espace mesur\'e. La \emph{compl\'etion} de $\mu$
est la mesure $\overline{\mu}$ d\'efinie sur la $\sigma$-alg\`ebre~:
\[
\overline{\mathcal{A}} = \{A \cup N : A \in \mathcal{A},\, N \subset M \text{ pour
un } M \in \mathcal{A} \text{ avec } \mu(M) = 0\}
\]
par $\overline{\mu}(A \cup N) = \mu(A)$.
\end{definition}

\begin{proposition}
$\overline{\mathcal{A}}$ est une $\sigma$-alg\`ebre et $\overline{\mu}$ est une
mesure compl\`ete qui \'etend $\mu$.
\end{proposition}

\begin{proof}
V\'erifions que $\overline{\mathcal{A}}$ est une $\sigma$-alg\`ebre.

\emph{Stabilit\'e par compl\'ementation}~: si $E = A \cup N$ avec $N \subset M$,
$\mu(M) = 0$, alors $E^c = A^c \cap N^c = (A^c \cap M^c) \cup (A^c \cap M
\setminus N)$. Or $A^c \cap M^c \in \mathcal{A}$ et $A^c \cap M \setminus N
\subset M$, donc $E^c \in \overline{\mathcal{A}}$.

\emph{Stabilit\'e par r\'eunion d\'enombrable}~: si $E_n = A_n \cup N_n$ avec
$N_n \subset M_n$, $\mu(M_n) = 0$, alors $\bigcup_n E_n = \bigcup_n A_n \cup
\bigcup_n N_n$ avec $\bigcup_n N_n \subset \bigcup_n M_n$ et
$\mu(\bigcup_n M_n) \leq \sum_n \mu(M_n) = 0$.

La bonne d\'efinition de $\overline{\mu}$ (ind\'ependance de la d\'ecomposition)
se v\'erifie directement, et la $\sigma$-additivit\'e d\'ecoule de celle de $\mu$.
\end{proof}

\section{Mesure image}

\begin{definition}[Mesure image]
Soient $(X, \mathcal{A}, \mu)$ un espace mesur\'e et $f : X \to Y$ une application
mesurable vers $(Y, \mathcal{B})$. La \emph{mesure image} (ou \emph{push-forward})
de $\mu$ par $f$ est la mesure $f_*\mu$ sur $(Y, \mathcal{B})$ d\'efinie par~:
\[
(f_*\mu)(B) = \mu(f^{-1}(B)), \quad B \in \mathcal{B}.
\]
\end{definition}

\begin{proposition}
$f_*\mu$ est bien une mesure sur $(Y, \mathcal{B})$.
\end{proposition}

\begin{proof}
$(f_*\mu)(\varnothing) = \mu(f^{-1}(\varnothing)) = \mu(\varnothing) = 0$. Pour
$(B_n)$ disjointe dans $\mathcal{B}$, $(f^{-1}(B_n))$ est disjointe dans
$\mathcal{A}$ et $f^{-1}(\bigsqcup_n B_n) = \bigsqcup_n f^{-1}(B_n)$, donc
$(f_*\mu)(\bigsqcup_n B_n) = \mu(\bigsqcup_n f^{-1}(B_n)) = \sum_n
\mu(f^{-1}(B_n)) = \sum_n (f_*\mu)(B_n)$.
\end{proof}

\section{Exercices}

\begin{exercice}
Soit $(X, \mathcal{A}, \mu)$ un espace mesur\'e et $(A_n) \subset \mathcal{A}$.
Montrer que
$\mu(\liminf_n A_n) \leq \liminf_n \mu(A_n)$
(analogue du lemme de Fatou pour les ensembles).
\end{exercice}

\begin{exercice}
Soit $\mu$ une mesure sur $(\R, \mathcal{B}(\R))$ telle que $\mu(K) < \infty$ pour
tout compact $K$. Montrer que $\mu$ est $\sigma$-finie.
\end{exercice}

\begin{exercice}
Montrer que la mesure de comptage sur $(\R, \mathcal{P}(\R))$ n'est pas $\sigma$-finie.
\end{exercice}

\begin{exercice}
Soit $(X, \mathcal{A}, \mu)$ un espace mesur\'e avec $\mu(X) < \infty$ et
$(A_n) \subset \mathcal{A}$ une suite telle que
$\mu(A_n) \geq \alpha > 0$ pour tout $n$. Montrer qu'il existe $x \in X$
appartenant \`a une infinit\'e de $A_n$.
\end{exercice}

\begin{exercice}
Soit $\mu$ une mesure finie sur $(\N, \mathcal{P}(\N))$. Montrer que $\mu$ est
d\'etermin\'ee par les valeurs $\mu(\{n\})$, $n \in \N$, et que
$\mu = \sum_{n=0}^{\infty} \mu(\{n\}) \delta_n$.
\end{exercice}

\begin{exercice}[Mesure de Lebesgue--Stieltjes]
Soit $F : \R \to \R$ une fonction croissante et continue \`a droite. On d\'efinit
$\mu_0((a, b]) = F(b) - F(a)$ sur l'alg\`ebre $\mathcal{A}_0$ des r\'eunions finies
disjointes d'intervalles semi-ouverts $(a, b]$.
\begin{enumerate}
    \item Montrer que $\mu_0$ est une pr\'e-mesure sur $\mathcal{A}_0$.
    \item En d\'eduire l'existence d'une unique mesure $\mu_F$ sur
    $\mathcal{B}(\R)$ telle que $\mu_F((a, b]) = F(b) - F(a)$.
\end{enumerate}
\end{exercice}
